\scanpage{75}
3) \uline{Essayer de trouver des définitions raisonnables en car.}2, \uline{ou de préférence indépendantes de la caractéristique}. Une théorie vraiment satisfaisante devrait marcher sur tout préschéma $S$, et réciproquement. Je suis loin d'avoir des idées claires sur la question, et me borne à des remarques vagues:

a) Lorsqu'on est sur un corps de caractéristique \deleted{2}\inserted{quelconque}, on est sans doute obligé de tenir compte des formes quadratiques dégénérées, et de définir un \deleted{\unsure{r}}anneau $Q'(k)$ comme défini par les classes de formes quadratiques, dégénérées ou non. \handwritten{[}Dans le cas d'un corps de \deleted{\illegible{1}} caractéristique $\handwritten{\ne}2$, ce serait $Q(k)+\underline{Z}$, où la \deleted{deuxième facteur $\underline{Z}$ est un} multiplication est donnée par $(\unsure{g|f},m).(g,n)=(fg,\deg(f)n+\deg(g)m+mn)$, le facteur en rabiot $\underline{Z}$ correspondant aux formes nulles, sur des vectoriels de rang quelconque\handwritten{]}.

b) En caractéristique quelconque, garder à l'esprit la distinction entre le shhéma en groupes $(\underline{Z}/2\underline{Z})_S$ et le shhéma en groupes $(\handwritten{\mu_2})_S$ sur $S$, le dernier étant le noyau de la puissance deuxième dans le groupe multiplicatif $\underline{G}_{m\,S}$. Si $\unsure{S|A}$ est local d'anneau $A$, alors la suite exacte de cohomologie donne $H^{\handwritten{1}}(S,\handwritten{\mu_2})=A^{\handwritten{*}}/A^{\handwritten{*2}}$, $H^{\handwritten{2}}(S,\handwritten{\mu_2})={}_{\handwritten{2}}H^{\handwritten{2}}(A,\underline{G}_m)$, éléments annulés par 2 dans le groupe de Brauer, par contre $H^{\handwritten{1}}(S,\underline{Z}/2\underline{Z})=A/\handwritten{\wp}A$ lorsque $2A=0$. Il faut donc s'attendre à trouver des classes de Stiefel-Whitney de type mixte, faisant à la fois intervenir les deux groupes voisins de coéfficients. Cette remarque se précise quand on regarde les schémas en groupes orthogonal\handwritten{s}, sur $\underline{Z}$ disons: en dimension paire, c'est l'extensions de $\underline{Z}/2\underline{Z}$ par un schéma en groupes à fibres connexes (le schéma spécial-orthogonal), par exemple en dim 2 c'est \deleted{l'extension} un produit semi-direct de $\underline{Z}/2\underline{Z}$ par $\underline{G}_m$ sur lequel il opère par symétrie; en dimension impaire, c'est une extension de $\handwritten{\mu_2}$ par le groupe spécial-orthogonal, par exemple en dimension 1, c'est $\handwritten{\mu_2}$. D'où suivant la parité, comme premier invariant ``discret'' d'un fibré principal orthhogonal, une 1-classe de groupe $\underline{Z}/2\underline{Z}$,ou $\handwritten{\mu_2}$.
% archive sheet number: 74
